\documentclass{handout}
\usepackage{handoutSetup}\usepackage{handoutShortcuts} % 
\usepackage[noheads,nolists,nomarkers]{endfloat}

\begin{document}
\handoutHeader

\begin{verbatimwrite}{\jobname.title}
Investment Notes
\end{verbatimwrite}

\handoutNameMake

\section{Broad Comments}

\subsection{Why Study Investment?}

\begin{enumerate}
\item Although investment is much smaller as a fraction of GDP than consumption, investment is 
much more {\it variable} than consumption.  So {\it fluctuations} in investment spending account
for  a large proportion of business-cycle frequency fluctuations in GDP.  (See the figure below for one way of measuring this).  



\item In the long run, the average magnitude of investment spending determines the size of the
capital stock.  If we think accounting for capital is an important component of explaining 
macroeconomic outcomes, then we need an understanding of investment
\end{enumerate}

\subsection{History of Thought}

\begin{itemize}
\item Keynes: $I$ determined by ``animal spirits'' which fluctuated strongly
\item \cite{samuelsonI}/Hansen/Harrod: Keynesian `Multiplier/Accelerator' model of investment
\item \cite{hall&jorgenson:i} neoclassical model with no costs of adjustment 
\item \cite{tobinsQ}'s $q$ model
\item \cite{abel:q}-\cite{hayashi:q} `marginal $\q$' model with smooth convex costs of adjustment
\item Models with nonconvex costs of adjustment
\item Models with capital market imperfections
\end{itemize}

\section{The Samuelson Multiplier-Accelerator Model}

From today's perspective, the distinctive feature of Samuelson's model of investment is its emphasis
on the proposition that levels of investment were determined by impulses from outside the market for
investment.  While the model is highly sophisticated in many ways, the essence of it for our purposes
can be captured by a very simple equation:

\begin{equation}\begin{gathered}\begin{aligned}
I_{t} & =  \alpha Y_{t} + \gamma \rfree_{t} \label{eq:samuelsonI}
\end{aligned}\end{gathered}\end{equation}

When paired with equations determining consumption and income, and with some lags added, the model
generated the classic Keynesian multiplier-accelerator framework.  But the key feature of the investment 
component of the model was that it assumed that the desired level of investment depended in a simple
linear way on the autonomous evolution of other variables: Output $Y$ and interest rates $\rfree$.  These variables
were in turn affected by $I$, and the model was capable of producing rich dynamics as a result.  
But for the purpose of explaining investment itself, the model's key implication could be boiled down to
\eqref{eq:samuelsonI}

In purely statistical terms, an equation like \eqref{eq:samuelsonI}
performs remarkably well in `explaining' investment spending.  But
this is not very satisfactory because income, investment, and interest
rates are presumably all determined by deeper underlying features of
the environment, and so it is highly problematic to interpret this
equation as instructive about the fundamental determinants of
investment.

\section{Hall and Jorgensen}

\cite{hall&jorgenson:i} made an important step forward by constructing a model in which an 
optimizing firm chooses the level of its capital stock with reference to fundamental features 
of the economic environment like its production function, depreciation rates, tax 
considerations, and a comparison of interest rates to the productivity of the firm's available 
uses of capital.  

\handoutI{HallJorgensen} presents a stripped-down version of this model.  

An important objection to this framework, however, was that it assumes that firms have the ability
to instantly and costlessly move their capital stocks to the level that would be justified by the
prevailing economic environment.  If something happens to make the firm wish to be larger, it can,
say, quadruple its capital stock overnight, paying no cost of adjustment for such a radical change.

\section{$\q$ Models}

A seminal paper by \cite{tobinsQ} provided the intuition for how to move forward.  Tobin defined
a variable

\begin{equation}\begin{gathered}\begin{aligned}
Q & =  \left(\frac{\text{Stock Mkt Val of Firm}}{\text{Replacement Cost of K}}\right)
\end{aligned}\end{gathered}\end{equation}
and proposed that investment should obey the rule
\begin{equation}
  \label{eq:1}
  \begin{cases}
   i_{t} > 0 & \text{if $Q_{t} > 1$}
\\ i_{t} < 0 & \text{if $Q_{t} < 1$}
\end{cases}
\end{equation}

That is, firms which have a value greater than what it would cost to reproduce
their capital should be growing, while firms which are not worth what it would cost
to reproduce them should be shrinking.

Tobin's $Q$ is now one of the basic tools of financial market analysis; it is used regularly
as one of several tools designed to try to assess firms' prospects.

Tobin's paper inspired a large amount of empirical and theoretical work, 
culminating in two classic papers, by \cite{abel:q} and \cite{hayashi:q},
which put the theory on a rigorous mathematical basis.  The canonical
model that emerged from this literature is summarized in \handoutI{qModel}.

\begin{comment} % Removed because it is a problem on the problem set!
\section{The Entrepreneurial Firm Under Perfect Foresight}

One unsatisfactory feature of the $\q$ model is that the motivations
of the firm's owners are not very clearly delineated.  Shareholders
are assumed to be risk-neutral return maximizers, and the firm chooses
its investment policies accordingly, but the model has no real
implications for the firm's dividend policy or holdings of cash.
(This is an example of the famous \cite{mmTheorem} theorem
derived under perfect capital markets).  Yet there are clearly
systematic patterns to the choices firms make about dividend payouts.  
In fact, much of what is studied in corporate finance is how to 
manage such payouts.  

One way to produce a well defined dividend payout policy is to consider
the firm as being wholly owned by a utility-maximizing entrepreneur whose
consumption is determined by the dividend payout from the firm.  That 
approach is explored in \handoutI{EntrepreneurPF}.  The surprising 
upshot is that, in the perfect foresight case, the behavior of the firm
owned by an entrepreneur is indistinguishable from the behavior of a ``publicly
traded'' firm, while the consumption behavior of the entrepreneur is 
indistinguishable from the behavior of a consumer who does not own any
business enterprises.  (This conclusion relies importantly on the assumption
of perfect foresight.)
\end{comment}

\section{Imperfect Capital Markets}

One of the foundations of the $\q$ model as developed by Abel and Hayashi 
is its reliance on an assumption that capital markets are perfect, in the sense
that in these models investment is pursued right up to the point where the 
marginal value of an additional unit of investment has been driven down 
to the exact value of the riskless interest rate (after accounting appropriately
for tax features, depreciation, etc), which reflects the optimal
quantity of investment.  

The key reason that the assumption of perfect capital markets is so useful
is that it implies that it is not necessary to keep track of the firm's 
cash management or financial activities.  

Useful though it may be, the assumption of perfect capital markets is clearly
very far from reality.  A vast literature has explored what happens when 
a wide variety of imperfections are introduced.  The handout \href{http://www.econ2.jhu.edu/people/ccarroll/Choice/LectureNotes/investment/DynanImperfCapMkts.pdf}{DynanImperfCapMkts.pdf} summarizes one
particular model of capital market imperfections that produces implications
that are similar to those that emerge from many other models.  

The implication of capital market imperfections that has been most 
vigorously explored is that the amount of investment may be affected
by a firm's financial position.  The crux of that literature is 
summarized in \handoutI{iAndCashFlow}.

% Figure from Torsten Slok at DB, in presentation for 2011/10, email from torsten.slok@db.com, Slok_Presentation_Oct.pdf 
\begin{figure}
\caption{Investment Drives the Cycle}
\begin{center}
\includegraphics[width=6in]{../Figures/InvestmentDrivesCycle} 
\end{center}
\end{figure}


\input handoutBibMake.tex

\end{document}

